\documentclass[11pt,a4paper]{article}

% ── Codificación, fuentes y lenguaje ───────────────────────────────────────────
\usepackage[utf8]{inputenc}
\usepackage[T1]{fontenc}
\usepackage[default,scale=0.95]{sourcesanspro}
\renewcommand{\familydefault}{\sfdefault}
\usepackage[spanish,es-noshorthands,es-tabla]{babel}

% ── Geometría y espaciado ─────────────────────────────────────────────────────
\usepackage[top=2.2cm, bottom=2.5cm, left=2.2cm, right=2.2cm, headheight=15pt]{geometry}
\usepackage{setspace}
\setstretch{1.18}
\usepackage{parskip}

% ── Matemáticas y unidades ────────────────────────────────────────────────────
\usepackage{amsmath}
\usepackage{amssymb}
\usepackage{siunitx}

% ── Circuitos ─────────────────────────────────────────────────────────────────
\usepackage[european, straightvoltages]{circuitikz}

% ── Tablas ────────────────────────────────────────────────────────────────────
\usepackage{booktabs}
\usepackage{tabularx}
\usepackage{array}
\usepackage{longtable}
\usepackage{colortbl}
\usepackage{enumitem}

% ── Colores, cajas e iconos ───────────────────────────────────────────────────
\usepackage[dvipsnames]{xcolor}
\usepackage{tcolorbox}
\tcbuselibrary{skins, breakable}
\usepackage{fontawesome5}
\usepackage{graphicx}

% ── Listados de código (dark-mode infográfico) ────────────────────────────────
%   Estilo base: fondo oscuro con números de línea. Los bloques lstlisting se
%   envuelven automáticamente en una caja tcolorbox redondeada.
\usepackage{listings}

% ── Secciones con barra de color (infografía) ────────────────────────────────
\usepackage{titlesec}
\titleformat{\section}
  {\normalfont\LARGE\bfseries\color{azulNoche}}
  {\colorbox{cyanNeon}{\color{fondoOscuro}\thesection}}{0.7em}{}
  [\vspace{2pt}\color{cyanNeon}\rule{\linewidth}{1.8pt}]
\titlespacing*{\section}{0pt}{24pt}{10pt}
\titleformat{\subsection}
  {\normalfont\large\bfseries\color{azulNoche}}
  {\textcolor{cyanNeon}{\thesubsection}}{0.7em}{}
  [\vspace{1pt}\color{grisLinea}\rule{\linewidth}{0.6pt}]
\titlespacing*{\subsection}{0pt}{14pt}{6pt}
\titleformat{\subsubsection}
  {\normalfont\normalsize\bfseries\color{azulNoche}}
  {\textcolor{cyanNeon}{\thesubsubsection}}{0.7em}{}
\titlespacing*{\subsubsection}{0pt}{10pt}{4pt}

% ── Cabeceras y pies de página ────────────────────────────────────────────────
\usepackage{fancyhdr}
\usepackage{lastpage}
\pagestyle{fancy}
\fancyhf{}
\renewcommand{\headrulewidth}{0pt}
\renewcommand{\footrulewidth}{0pt}
\fancyfoot[C]{%
  \begin{minipage}{\linewidth}
    \vspace{2pt}
    \color{cyanNeon}\rule{\linewidth}{1.2pt}\\[2pt]
    \centering\color{grisTexto}\footnotesize
    Página \thepage\ de \pageref{LastPage}
  \end{minipage}%
}

% ── Hiperenlaces ──────────────────────────────────────────────────────────────
\usepackage[colorlinks=true, linkcolor=azulNoche, urlcolor=cyanNeon]{hyperref}
\sloppy
\emergencystretch 3em

% ── Paleta de colores — tecnológico dark-mode ──────────────────────────────────
\definecolor{fondoOscuro}{HTML}{0D1B2A}%             Dark navy — fondo banda título
\definecolor{verdeTurquesa}{HTML}{004D40}%           Verde turquesa — bandas de marca
\definecolor{azulNoche}{HTML}{1B3A6B}%               Midnight blue — secciones, marcos
\definecolor{azulMedio}{HTML}{2A5298}%               Azul medio — variante de acento
\definecolor{cyanNeon}{HTML}{00C8E0}%                Cyan eléctrico — acentos primarios
\definecolor{verdeSignal}{HTML}{00B887}%             Verde señal — fortalezas, éxito
\definecolor{naranjaVivo}{HTML}{FF7043}%             Naranja vivo — mejoras, advertencias
\definecolor{rojoAlerta}{HTML}{E53935}%              Rojo alerta — errores
\definecolor{grisPapel}{HTML}{F4F6FA}%               Fondo tarjetas claras
\definecolor{grisLinea}{HTML}{C8D0E0}%               Bordes sutiles
\definecolor{grisTexto}{HTML}{6B7A99}%               Texto secundario
\definecolor{amarilloNota}{HTML}{FFB300}%            Notas / advertencias doradas
\definecolor{codigoFondo}{HTML}{1E2D3D}%             Fondo del bloque de código
\definecolor{codigoComentario}{HTML}{7A8FA6}%        Comentarios
\definecolor{codigoCadena}{HTML}{FF9D5C}%            Cadenas / strings
\definecolor{codigoKeyword}{HTML}{00C8E0}%           Palabras clave
\definecolor{codigoNumero}{HTML}{00E5A0}%            Números

% ── Banda de título: portada infográfica ──────────────────────────────────────
%   Diseño en 2 capas: banda de color (por defecto fondo oscuro) + franja de
%   acento cyan abajo. Uso: \bandaTitulo[color]{título}{subtítulo}.
%   El icono se puede cambiar con \renewcommand{\iconoBanda}{...} (FontAwesome).
\newcommand{\iconoBanda}{microchip}
\newcommand{\bandaTitulo}[3][fondoOscuro]{%
  \begin{tcolorbox}[
    enhanced,
    colback=#1, colframe=#1,
    boxrule=0pt, arc=8pt, outer arc=8pt,
    width=\linewidth,
    top=18pt, bottom=6pt, left=14pt, right=14pt,
    borderline south={3.5pt}{0pt}{cyanNeon},
  ]
    \begin{minipage}[c]{0.07\linewidth}
      \centering
      \textcolor{cyanNeon}{\fontsize{30}{30}\selectfont\faIcon{\iconoBanda}}
    \end{minipage}%
    \hspace{8pt}%
    \begin{minipage}[c]{0.88\linewidth}
      {\fontsize{20}{24}\selectfont\bfseries\color{white}#2}\par\vspace{5pt}%
      \textcolor{cyanNeon!80!white}{\small\faIcon{angle-right}~#3}%
    \end{minipage}
    \vspace{4pt}
  \end{tcolorbox}%
  \vspace{8pt}%
}

% ── Tarjetas de datos (infografía) ────────────────────────────────────────────
\tcbset{
  tarjetaDato/.style={
    enhanced, breakable,
    arc=6pt, outer arc=6pt,
    colback=grisPapel, colframe=azulNoche,
    boxrule=1pt,
    fonttitle=\bfseries\small, coltitle=white,
    attach boxed title to top left={yshift=-3mm, xshift=6mm},
    boxed title style={
      colback=azulNoche, colframe=azulNoche,
      arc=4pt, boxrule=0pt,
    },
    drop fuzzy shadow=azulNoche!30!white,
    top=8pt, bottom=8pt, left=10pt, right=10pt,
  },
  tarjetaIcono/.style={
    enhanced, breakable,
    arc=6pt, outer arc=6pt,
    colback=white, colframe=grisLinea, boxrule=0.8pt,
    fonttitle=\bfseries\small, coltitle=azulNoche,
    borderline west={3pt}{0pt}{cyanNeon},
    drop fuzzy shadow=azulNoche!25!white,
    top=6pt, bottom=6pt, left=10pt, right=10pt,
  },
  cajaTitulo/.style={
    enhanced, breakable,
    arc=6pt, outer arc=6pt,
    colback=azulNoche!10!grisPapel, colframe=azulNoche,
    boxrule=1pt,
    fonttitle=\bfseries, coltitle=white,
    attach boxed title to top left={yshift=-3mm, xshift=6mm},
    boxed title style={
      colback=azulNoche, colframe=azulNoche,
      arc=4pt, boxrule=0pt,
    },
    drop fuzzy shadow=azulNoche!25!white,
    top=8pt, bottom=8pt, left=10pt, right=10pt,
  },
  cajaContenido/.style={
    enhanced, breakable,
    arc=5pt, outer arc=5pt,
    colback=grisPapel, colframe=grisLinea,
    boxrule=0.8pt,
    fonttitle=\bfseries\small, coltitle=azulNoche,
    top=6pt, bottom=6pt, left=10pt, right=10pt,
  },
  cajaFortaleza/.style={
    enhanced, breakable,
    arc=6pt, outer arc=6pt,
    colback=verdeSignal!8!white, colframe=verdeSignal,
    boxrule=1pt,
    borderline west={4pt}{0pt}{verdeSignal},
    fonttitle=\bfseries\small, coltitle=verdeSignal!70!black,
    drop fuzzy shadow=verdeSignal!20!white,
    top=6pt, bottom=6pt, left=10pt, right=10pt,
  },
  cajaMejora/.style={
    enhanced, breakable,
    arc=6pt, outer arc=6pt,
    colback=naranjaVivo!8!white, colframe=naranjaVivo,
    boxrule=1pt,
    borderline west={4pt}{0pt}{naranjaVivo},
    fonttitle=\bfseries\small, coltitle=naranjaVivo!80!black,
    drop fuzzy shadow=naranjaVivo!20!white,
    top=6pt, bottom=6pt, left=10pt, right=10pt,
  },
  cajaRecomendacion/.style={
    enhanced, breakable,
    arc=6pt, outer arc=6pt,
    colback=cyanNeon!6!white, colframe=cyanNeon!60!azulNoche,
    boxrule=1pt,
    borderline west={4pt}{0pt}{cyanNeon},
    fonttitle=\bfseries\small, coltitle=azulNoche,
    drop fuzzy shadow=azulNoche!20!white,
    top=6pt, bottom=6pt, left=10pt, right=10pt,
  },
}

% ── Ícono + texto (encabezados de sección y elementos) ────────────────────────
\newcommand{\iconotexto}[2]{\textcolor{cyanNeon}{\faIcon{#1}}~#2}

% ── Listados de código: estilo dark + auto-envoltura ─────────────────────────
\lstdefinestyle{estiloCodigo}{
    backgroundcolor=\color{codigoFondo},
    basicstyle=\ttfamily\small\color{white},
    commentstyle=\color{codigoComentario}\itshape,
    stringstyle=\color{codigoCadena},
    keywordstyle=\color{codigoKeyword}\bfseries,
    numberstyle=\tiny\color{codigoComentario},
    numbers=left,
    numbersep=10pt,
    showstringspaces=false,
    breaklines=true,
    breakatwhitespace=false,
    tabsize=4,
    frame=none,
    captionpos=b,
    aboveskip=6pt,
    belowskip=4pt,
    xleftmargin=14pt,
    framexleftmargin=14pt,
    literate=
      {á}{{\'a}}1 {é}{{\'e}}1 {í}{{\'i}}1 {ó}{{\'o}}1 {ú}{{\'u}}1
      {Á}{{\'A}}1 {É}{{\'E}}1 {Í}{{\'I}}1 {Ó}{{\'O}}1 {Ú}{{\'U}}1
      {ñ}{{\~n}}1 {Ñ}{{\~N}}1 {ü}{{\"u}}1 {Ü}{{\"U}}1
      {¿}{{?`}}1 {¡}{{!`}}1
}
\lstdefinestyle{estiloPython}{
    style=estiloCodigo,
    language=Python,
}
\lstset{style=estiloCodigo}
\BeforeBeginEnvironment{lstlisting}{%
  \begin{tcolorbox}[
    enhanced, breakable,
    arc=6pt, outer arc=6pt,
    colback=codigoFondo, colframe=azulNoche,
    boxrule=1pt,
    drop fuzzy shadow=azulNoche!25!white,
    top=2pt, bottom=2pt, left=0pt, right=0pt,
  ]%
}
\AfterEndEnvironment{lstlisting}{\end{tcolorbox}}

% ── Cajas de contenido temáticas ──────────────────────────────────────────────
\newtcolorbox{cajaRecuerda}{
    enhanced, breakable,
    arc=6pt, outer arc=6pt,
    colback=azulNoche!10!grisPapel,
    colframe=azulNoche, boxrule=1pt,
    borderline west={4pt}{0pt}{cyanNeon},
    fonttitle=\bfseries\small\color{white},
    title={\faIcon{info-circle}~Recuerda},
    attach boxed title to top left={yshift=-3mm, xshift=6mm},
    boxed title style={colback=azulNoche, arc=4pt, boxrule=0pt},
    drop fuzzy shadow=azulNoche!25!white,
    left=10pt, right=10pt, top=8pt, bottom=8pt,
}

\newtcolorbox{cajaConcepto}{
    enhanced, breakable,
    arc=6pt, outer arc=6pt,
    colback=verdeSignal!8!white,
    colframe=verdeSignal, boxrule=1pt,
    borderline west={4pt}{0pt}{verdeSignal},
    fonttitle=\bfseries\small\color{white},
    title={\faIcon{lightbulb}~Concepto clave},
    attach boxed title to top left={yshift=-3mm, xshift=6mm},
    boxed title style={colback=verdeSignal!80!black, arc=4pt, boxrule=0pt},
    drop fuzzy shadow=verdeSignal!20!white,
    left=10pt, right=10pt, top=8pt, bottom=8pt,
}

\newtcolorbox{cajaEjemplo}{
    enhanced, breakable,
    arc=6pt, outer arc=6pt,
    colback=cyanNeon!5!white,
    colframe=cyanNeon!70!azulNoche, boxrule=1pt,
    borderline west={4pt}{0pt}{cyanNeon},
    fonttitle=\bfseries\small\color{fondoOscuro},
    title={\faIcon{code}~Ejemplo},
    attach boxed title to top left={yshift=-3mm, xshift=6mm},
    boxed title style={colback=cyanNeon, arc=4pt, boxrule=0pt},
    drop fuzzy shadow=azulNoche!20!white,
    left=10pt, right=10pt, top=8pt, bottom=8pt,
}

\newtcolorbox{cajaEjercicio}{
    enhanced, breakable,
    arc=6pt, outer arc=6pt,
    colback=azulMedio!7!white,
    colframe=azulMedio, boxrule=1pt,
    borderline west={4pt}{0pt}{azulMedio},
    fonttitle=\bfseries\small\color{white},
    title={\faIcon{pencil-alt}~Ejercicio},
    attach boxed title to top left={yshift=-3mm, xshift=6mm},
    boxed title style={colback=azulMedio, arc=4pt, boxrule=0pt},
    drop fuzzy shadow=azulNoche!20!white,
    left=10pt, right=10pt, top=8pt, bottom=8pt,
}

\newtcolorbox{cajaReto}{
    enhanced, breakable,
    arc=6pt, outer arc=6pt,
    colback=naranjaVivo!6!white,
    colframe=naranjaVivo, boxrule=1pt,
    borderline west={4pt}{0pt}{naranjaVivo},
    fonttitle=\bfseries\small\color{white},
    title={\faIcon{fire}~Reto},
    attach boxed title to top left={yshift=-3mm, xshift=6mm},
    boxed title style={colback=naranjaVivo, arc=4pt, boxrule=0pt},
    drop fuzzy shadow=naranjaVivo!20!white,
    left=10pt, right=10pt, top=8pt, bottom=8pt,
}

% ── Indicadores de dificultad ─────────────────────────────────────────────────
\newcommand{\facil}{\textcolor{verdeSignal}{\faIcon{circle}\,\textbf{Fácil}}}
\newcommand{\intermedio}{\textcolor{naranjaVivo}{\faIcon{circle}\,\textbf{Intermedio}}}
\newcommand{\dificil}{\textcolor{rojoAlerta}{\faIcon{circle}\,\textbf{Difícil}}}

% ─────────────────────────────────────────────────────────────────────────────
\renewcommand{\iconoBanda}{atom}
\fancyhead[L]{\footnotesize\color{grisTexto}\faIcon{atom}~LABORATORIO\_II\_FISICA}
\fancyhead[R]{\footnotesize\color{grisTexto}Práctica 06}
\hypersetup{pdftitle={Laboratorio II de Física},pdfauthor={Prof. César Rodríguez}
  pdfsubject={LABORATORIO\_II\_FISICA}}

\begin{document}
\bandaTitulo[fondoOscuro]{Práctica 06 --- Determinación de reactancias, de C y de L en régimen senoidal}{Semana 6 --- Fase 2: Elementos no lineales y régimen senoidal}

\begin{center}\small\color{grisTexto}Prof. César Rodríguez \textbar{} Universidad de Oriente, Núcleo de Sucre, Departamento de Física \textbar{} 2 de noviembre de 2026\end{center}

\begin{tabularx}{\linewidth}{@{}>{\raggedright\arraybackslash}p{4.6cm}X@{}}
\toprule
\textbf{Duración} & 3 horas de laboratorio presencial más 2 horas de análisis \\
\textbf{Modalidad} & 4 grupos con roles rotativos \\
\textbf{Entrega} & GitHub del grupo, carpeta semana\_06, antes del corte semanal \\
\textbf{Instrumento} & Práctica formativa con informe \\
\bottomrule
\end{tabularx}

\section{\iconotexto{bullseye}{Objetivo}}

Estudiar el comportamiento de los circuitos RC y RL en régimen senoidal y usar esos circuitos para determinar la capacitancia de un capacitor y la inductancia de una bobina.

\section{\iconotexto{clipboard-list}{Contenidos que se aplican}}

\begin{itemize}[leftmargin=1.4em,itemsep=2pt,topsep=3pt]
  \item Señal senoidal, amplitud, frecuencia angular, fase y valor eficaz.
  \item Concepto de fasor e impedancia compleja.
  \item Reactancia inductiva y capacitiva; resistencia como caso particular.
  \item Definición de reactancia como cociente entre tensión y corriente.
  \item Dependencia de la reactancia con la frecuencia angular.
  \item Impedancia equivalente serie y paralelo.
  \item Resistencia serie de un inductor real y capacitancia distribuida de una bobina.
\end{itemize}

\section{\iconotexto{book-open}{Fundamento teórico}}

En régimen senoidal permanente cada elemento se describe por su impedancia
compleja:
\[
  Z_R = R \qquad
  Z_L = j\omega L \qquad
  Z_C = \frac{1}{j\omega C} = -\frac{j}{\omega C}
\]
con $\omega = 2\pi f$. Los módulos de las reactancias son
\[
  X_L = \omega L \qquad\qquad X_C = \frac{1}{\omega C}
\]
La corriente se adelanta 90 grados a la tensión en el inductor y se atrasa
90 grados en el capacitor; en el resistor están en fase. Para un circuito
serie la impedancia total y el ángulo de fase son
\[
  Z = \sqrt{R^2 + \left(X_L - X_C\right)^2}
  \qquad\qquad
  \varphi = \arctan\!\frac{X_L - X_C}{R}
\]
Un inductor real incluye la resistencia de su devanado, por lo que su
impedancia medida no es puramente imaginaria:
$Z_L = \sqrt{R_s^2 + X_L^2}$. Con osciloscopio de dos canales, la diferencia
de fase entre las señales se obtiene del desplazamiento temporal
$\Delta t$ respecto del periodo $T$ mediante $\varphi = 360^\circ\,\Delta t/T$.

\begin{center}
\begin{circuitikz}
  \draw (0,0) to[V, l=$v_g$] (0,3) to[R, l=$R$] (2.4,3)
        to[L, l=$L$] (4.8,3) to[C, l=$C$] (7.2,3) -- (7.2,0) -- (0,0);
  \draw (2.4,3) -- (2.4,4.2) to[voltmeter, l=$v_R$] (0.6,4.2);
\end{circuitikz}
\par\smallskip{\small\itshape Circuito serie alimentado por el generador; la tensión sobre R es proporcional a la corriente}
\end{center}

\section{\iconotexto{box-open}{Materiales y equipos}}

\begin{itemize}[leftmargin=1.4em,itemsep=2pt,topsep=3pt]
  \item Un generador de funciones senoidal
  \item Un resistor variable del tipo décadas
  \item Un capacitor del inventario del laboratorio
  \item Una bobina del inventario del laboratorio
  \item Dos multimedidores Keithley (175 o 175-A)
\end{itemize}

\section{\iconotexto{flask}{Escenario}}

Con un generador de funciones senoidal, una resistencia de décadas y dos Keithley, el grupo mide la reactancia de un capacitor y de una bobina del inventario del laboratorio, y después determina su capacitancia y su inductancia variando la frecuencia entre 200 y 1200 Hz. Los valores obtenidos se contrastan con los del apéndice del inventario.

\section{\iconotexto{list-ol}{Procedimiento}}

\begin{enumerate}[leftmargin=1.4em,itemsep=2pt,topsep=3pt]
  \item Armar el circuito RC serie y hacerlo revisar por el profesor antes de energizar.
  \item Fijar la frecuencia del generador alrededor de 200 Hz y un valor de resistencia.
  \item Tomar veinte pares de corriente $I$ y tensión $V_C$ moviendo el control de amplitud del generador.
  \item Armar el circuito RL serie y repetir los pasos anteriores tomando veinte pares de $I$ y $V_L$ a la misma frecuencia.
  \item Para determinar C, fijar un valor de R y variar la frecuencia veinte veces entre 200 y 1200 Hz, midiendo $V_C$ y $V_R$ en cada punto.
  \item Para determinar L, repetir el barrido de frecuencia midiendo $V_L$ y $V_R$ en cada punto.
  \item Calcular la corriente en cada punto como $I = V_R/R$ y anotar la incertidumbre de las dos lecturas.
\end{enumerate}

\section{\iconotexto{table}{Tabla de datos}}

\begin{center}
\begin{minipage}{\linewidth}
\centering
{\small
\textbf{Barrido de frecuencia para determinar C y L}\par\smallskip
\begin{tabular}{@{}ccccc@{}}
\toprule
$f$ (Hz) & $V_R$ (V) & $V_C$ (V) & $I$ (mA) & $X_C$ ($\Omega$) \\
\midrule
200 & $\ldots$ & $\ldots$ & $\ldots$ & $\ldots$ \\
300 & $\ldots$ & $\ldots$ & $\ldots$ & $\ldots$ \\
400 & $\ldots$ & $\ldots$ & $\ldots$ & $\ldots$ \\
500 & $\ldots$ & $\ldots$ & $\ldots$ & $\ldots$ \\
600 & $\ldots$ & $\ldots$ & $\ldots$ & $\ldots$ \\
800 & $\ldots$ & $\ldots$ & $\ldots$ & $\ldots$ \\
1000 & $\ldots$ & $\ldots$ & $\ldots$ & $\ldots$ \\
1200 & $\ldots$ & $\ldots$ & $\ldots$ & $\ldots$ \\
\bottomrule
\end{tabular}
}\par\smallskip{\footnotesize Los experimentos 1 y 2 usan frecuencia fija y amplitud variable, con su propia tabla de veinte filas.}
\end{minipage}
\end{center}

\section{\iconotexto{calculator}{Análisis y tratamiento de datos}}

\begin{enumerate}[leftmargin=1.4em,itemsep=2pt,topsep=3pt]
  \item Demostrar a partir de la definición de reactancia que $|X_C| = 1/(C\omega)$ y que $|X_L| = L\omega$.
  \item Graficar $V_C$ frente a $I$ con los veinte puntos de frecuencia fija y ajustar la recta por mínimos cuadrados.
  \item Obtener $X_C$ de la pendiente con su error y compararlo con el valor teórico calculado a partir de C.
  \item Repetir el mismo tratamiento con $V_L$ frente a $I$ para obtener $X_L$.
  \item Graficar $X_C$ frente a $1/\omega$ y $X_L$ frente a $\omega$, ajustar y obtener C y L de las pendientes, con sus errores.
  \item Calcular la desviación porcentual de C frente al valor de la tabla de capacitores del apéndice y de L frente a las tablas de inductancias.
\end{enumerate}

\section{\iconotexto{plus-circle}{Ampliación (opcional)}}

\begin{itemize}[leftmargin=1.4em,itemsep=2pt,topsep=3pt]
  \item Medir el desfase entre tensión y corriente con osciloscopio de dos canales y compararlo con el previsto por el modelo fasorial.
  \item Verificar el desfase con una figura de Lissajous y cuantificar la discrepancia entre ambos procedimientos.
\end{itemize}

\section{\iconotexto{box-open}{Entregable}}

Gráficas de $V_C$ y $V_L$ frente a $I$ con sus ajustes, valores de $X_C$ y $X_L$, y de C y L con su incertidumbre, comparados con el inventario del laboratorio.

\section{\iconotexto{ruler-combined}{Rúbrica de corrección}}

\begin{tabularx}{\linewidth}{@{}X r@{}}
\toprule
Criterio & Peso\\
\midrule
Fundamento teórico y modelo aplicado & 1,0 \\
Montaje, seguridad y manejo del instrumental & 1,5 \\
Toma de datos completa y trazable & 2,0 \\
Presupuesto de incertidumbres & 2,0 \\
Análisis, ajuste y gráficas & 2,0 \\
Conclusiones contrastadas con el modelo & 1,0 \\
Presentación y estructura del informe & 0,5 \\
\midrule
\textbf{Total} & \textbf{10,0}\\
\bottomrule
\end{tabularx}

\section{\iconotexto{code}{Contrato del script}}

\begin{cajaConcepto}Todo script entregado imprime \textbf{un solo objeto JSON} por la salida estándar y devuelve un diccionario con estas claves. \texttt{status} vale \texttt{ok} o \texttt{error}; si es \texttt{error}, \texttt{errores} explica el motivo y el programa sale con código distinto de cero.
\begin{lstlisting}[language=Python,style=estiloPython]
{
  "status": "ok",
  "practica": 1,
  "magnitud": "resistencia",
  "valor": 99.4,
  "unidad": "ohm",
  "incertidumbre": 0.6,
  "metodo": "ajuste_lineal_minimos_cuadrados",
  "n_puntos": 8,
  "residual_max": 0.012,
  "fuente_datos": "semana_01/datos/iv_resistor_220.csv",
  "errores": []
}

\end{lstlisting}
\end{cajaConcepto}

\begin{tcolorbox}[cajaMejora]\textbf{Antes de empezar:} Antes de energizar, contrasta el montaje con la figura y verifica rangos y polaridad de cada instrumento. Si un instrumento se calienta, huele o marca fuera de escala, corta la alimentación y avisa al docente: el montaje se revisa antes de continuar.
\end{tcolorbox}
\end{document}
